\documentclass[11pt,a4paper]{article}
\usepackage[utf8]{inputenc}
\usepackage[T1]{fontenc}
\usepackage[english]{babel}
\usepackage{amsmath}
\usepackage{amssymb}
\usepackage[margin=1.8cm]{geometry}
\usepackage{lmodern}
\usepackage{xcolor}
\usepackage{setspace}
\pagestyle{empty}
\setlength{\parindent}{0pt}
\setlength{\parskip}{0.8em}
\linespread{1.15}

\newcommand{\answernote}[1]{%
\vfill
\noindent\textcolor{gray}{\rule{\linewidth}{0.4pt}}

{\small\textcolor{gray}{Correct answer: #1}}
}

\begin{document}

{\LARGE\textbf{Question 1}}

For the given trajectory:

\[
\vec r(t) = \bigl(2t^2 - t, t^3 - 3t\bigr)
\]

the acceleration vector at time $t=2$ is:

\vspace{0.5em}

{\large
\textbf{A.} $(6, 12)$

\textbf{B.} $(4, 12)$

\textbf{C.} $(6, 6)$

\textbf{D.} $(8, 12)$
}

\answernote{B}
\newpage

{\LARGE\textbf{Question 2}}

For the given velocity and initial condition:

\[
\vec v(t) = \bigl(3t^2, 2t\bigr), \quad \vec r(0)=(0,0)
\]

the position at time $t=2$ is:

\vspace{0.5em}

{\large
\textbf{A.} $(8, 4)$

\textbf{B.} $(6, 4)$

\textbf{C.} $(8, 2)$

\textbf{D.} $(4, 2)$
}

\answernote{A}
\newpage

{\LARGE\textbf{Question 3}}

For the given acceleration and initial condition:

\[
\vec a(t) = \bigl(6t, -4\bigr), \quad \vec v(0)=(0,0)
\]

the velocity at time $t=2$ is:

\vspace{0.5em}

{\large
\textbf{A.} $(12, -8)$

\textbf{B.} $(6, -8)$

\textbf{C.} $(12, -4)$

\textbf{D.} $(6, -4)$
}

\answernote{A}
\newpage

{\LARGE\textbf{Question 4}}

For the given trajectory and the mass of the body:

\[
\vec r(t) = \bigl(t^3 - 2t, 2t^2 + t\bigr), \quad m = 2
\]

the force acting on the body at time $t=4$ is:

\vspace{0.5em}

{\large
\textbf{A.} $(48, 8)$

\textbf{B.} $(24, 4)$

\textbf{C.} $(48, 4)$

\textbf{D.} $(24, 8)$
}

\answernote{A}
\newpage

{\LARGE\textbf{Question 5}}

For the given trajectory and position-dependent force:

\[
x(t) = t^2, \quad t \in [0,1]
\]

\[
F(x) = 2x
\]

the work done by this force during the first second of motion is:

\vspace{0.5em}

{\large
\textbf{A.} $\frac{1}{2}$

\textbf{B.} $\frac{2}{3}$

\textbf{C.} $1$

\textbf{D.} $\frac{3}{2}$
}

\answernote{C}
\newpage

{\LARGE\textbf{Question 6}}

For the given trajectory:

\[
\vec r(t) = \bigl(\cos(t^2), \sin(t^2)\bigr), \qquad t \ge 0
\]

which of the following statements is true?

\vspace{0.5em}

{\large
\textbf{A.} The motion takes place along a circle, and the speed is constant

\textbf{B.} The motion takes place along a circle, and the speed increases linearly with time

\textbf{C.} The motion does not take place along a circle, and the speed is constant

\textbf{D.} The motion does not take place along a circle, and the speed increases linearly with time
}

\answernote{B}
\newpage

{\LARGE\textbf{Question 7}}

A tennis ball is dropped from the height of the 4th floor. Assume that one floor has a height of $3\ \text{m}$. After each bounce, the ball loses $\frac{1}{4}$ of its mechanical energy. To what maximum height will the ball rise after the third bounce?

\vspace{0.5em}

{\large
\textbf{A.} $12\cdot\left(\frac{3}{4}\right)^2$

\textbf{B.} $12\cdot\left(\frac{3}{4}\right)^3$

\textbf{C.} $12\cdot\left(\frac{1}{4}\right)^2$

\textbf{D.} $12\cdot\left(\frac{1}{4}\right)^3$
}

\answernote{B}
\newpage

{\LARGE\textbf{Question 8}}

For a harmonic oscillator with an initial maximum displacement $A$, which loses $\frac{1}{4}$ of its mechanical energy each time it passes through the equilibrium position, the maximum displacement reached \textbf{immediately after the second passage through the equilibrium position} is:

\vspace{0.5em}

{\large
\textbf{A.} $A\left(\frac{3}{4}\right)^2$

\textbf{B.} $A\sqrt{\frac{3}{4}}$

\textbf{C.} $\frac{3}{4}A$

\textbf{D.} $\frac{1}{2}A$
}

\answernote{C}
\newpage

{\LARGE\textbf{Question 9}}

For the given wave:

\[
y(x,t)=2\sin(\pi x-4\pi t)
\]

which of the following statements is true?

\vspace{0.5em}

{\large
\textbf{A.} The wave propagates with speed $4$, and the points $x=0$ and $x=2$ oscillate in the same phase

\textbf{B.} The wave propagates with speed $2$, and the points $x=0$ and $x=2$ oscillate in the same phase

\textbf{C.} The wave propagates with speed $2$, and the points $x=0$ and $x=2$ oscillate in opposite phase

\textbf{D.} The wave propagates with speed $4$, and the points $x=0$ and $x=2$ oscillate in opposite phase
}

\answernote{A}
\newpage

{\LARGE\textbf{Question 10}}

Two oscillations are given:

\[
y_1(t)=A\sin(2\pi f_1 t), \qquad y_2(t)=A\sin(2\pi f_2 t),
\]

where $f_1=10\ \text{Hz}$ and $f_2=12\ \text{Hz}$.

The beat frequency is:

\vspace{0.5em}

{\large
\textbf{A.} $1\ \text{Hz}$

\textbf{B.} $2\ \text{Hz}$

\textbf{C.} $11\ \text{Hz}$

\textbf{D.} $22\ \text{Hz}$
}

\answernote{B}

\end{document}
